\section{Mục tiêu đề tài}

Đề tài tập trung vào việc nghiên cứu, triển khai và đánh giá hiệu quả của các thuật toán tìm kiếm trong trí tuệ nhân tạo thông qua trò chơi Cờ vua. Mục tiêu chính là xây dựng một hệ thống AI có khả năng thi đấu ổn định và so sánh sự khác biệt giữa hai trường phái thuật toán phổ biến: duyệt cây có hệ thống và mô phỏng xác suất.

Các mục tiêu cụ thể bao gồm:
\begin{itemize}
    \item \textbf{Triển khai thuật toán Alpha-Beta Pruning:} Xây dựng công cụ tìm kiếm dựa trên thuật toán Minimax kết hợp cắt tỉa Alpha-Beta để tối ưu hóa không gian trạng thái. Đồng thời, thiết kế hàm đánh giá (Heuristic) hiệu quả dựa trên giá trị quân cờ (Material Score) và bảng điểm vị trí (Position Square Tables - PST).
    \item \textbf{Triển khai thuật toán Monte Carlo Tree Search (MCTS):} Xây dựng quy trình tìm kiếm dựa trên xác suất với đầy đủ 4 giai đoạn: Chọn lọc (Selection - sử dụng UCB1), Mở rộng (Expansion), Mô phỏng (Simulation) và Truyền ngược (Backpropagation). Đảm bảo AI có khả năng đưa ra nước đi hợp lý dựa trên thống kê thắng/thua.
    \item \textbf{Phát triển giao diện tương tác:} Xây dựng giao diện web (HTML/CSS/JS) trực quan, sinh động, cho phép người dùng thi đấu trực tiếp với AI và tùy chọn thuật toán đối đầu. Đảm bảo tính đồng bộ giữa logic xử lý cờ vua và hiển thị đồ họa.
    \item \textbf{Đánh giá và so sánh:} Phân tích ưu, nhược điểm của từng thuật toán về độ sâu tìm kiếm, thời gian phản hồi và chất lượng nước đi trong các tình huống thực tế, từ đó rút ra kết luận về tính ứng dụng của mỗi phương pháp.
\end{itemize}
